% Mathématiques 3e — V3 LaTeX
% Grandeurs composées et conversions

\section{Objectifs}
\begin{itemize}
\item Utiliser et interpréter des grandeurs quotients ou produits.
\item Convertir des unités composées avec cohérence.
\item Réinvestir les puissances de 10 dans les conversions de surfaces et volumes.
\item Résoudre des problèmes mêlant plusieurs grandeurs.
\end{itemize}

\section{Repères et vocabulaire}
Une grandeur composée associe plusieurs unités et décrit une relation physique ou pratique : kilomètres par heure, litres par minute, kilogrammes par mètre cube. En troisième, les conversions doivent être justifiées par le sens des unités et par les puissances de dix.

\section{1. Vitesse}
La vitesse moyenne est une distance divisée par une durée.

\[
v=\dfrac dt
\]

\begin{exemple}
Les unités de $d$ et $t$ déterminent l’unité de $v$ ; on ne mélange pas kilomètres et secondes sans conversion.
\end{exemple}

\section{2. Débit}
Un débit est une quantité transférée par unité de temps.

\[
q=\dfrac Vt
\]

\begin{exemple}
Un débit de $12\,\text{L/min}$ signifie que $12$ litres sont transférés en une minute dans le modèle constant.
\end{exemple}

\coursefigure{/images/mathematiques/3eme/grandeurs-composees-conversions-1.svg}{Représentation structurante pour le chapitre Grandeurs composées et conversions.}{Cette figure organise visuellement les relations utilisées dans les premières notions du chapitre.}

\section{3. Masse volumique}
La masse volumique relie masse et volume.

\[
\rho=\dfrac mV
\]

\begin{exemple}
La valeur numérique dépend des unités choisies : $1\,\text{g/cm}^3=1000\,\text{kg/m}^3$.
\end{exemple}

\section{4. Concentration}
Une concentration peut exprimer une masse de soluté par volume de solution.

\[
c=\dfrac mV
\]

\begin{exemple}
Il faut préciser l’unité, par exemple $\text{g/L}$, et distinguer concentration et masse volumique.
\end{exemple}

\section{5. Conversions linéaires}
Une conversion de longueur multiplie l’unité par une puissance de $10$.

\[
1\,\text{m}=10^2\,\text{cm}
\]

\begin{exemple}
Pour les surfaces et volumes, le facteur doit être élevé au carré ou au cube.
\end{exemple}

\coursefigure{/images/mathematiques/3eme/grandeurs-composees-conversions-2.svg}{Deuxième représentation pédagogique pour le chapitre Grandeurs composées et conversions.}{Cette représentation sert de support de lecture, de comparaison ou de contrôle dans les problèmes du chapitre.}

\section{6. Conversions de volumes}
Puisque $1\,\text{m}=10^2\,\text{cm}$, on obtient $1\,\text{m}^3=10^6\,\text{cm}^3$.

\[
(10^2)^3=10^6
\]

\begin{exemple}
Cette justification évite les conversions fausses par simple déplacement de virgule.
\end{exemple}

\section{7. Vitesse en unités différentes}
Pour convertir $90\,\text{km/h}$ en mètres par seconde, on convertit simultanément la distance et la durée.

\[
90\,\text{km/h}=25\,\text{m/s}
\]

\begin{exemple}
Le facteur usuel $3{,}6$ peut être retrouvé par $1000/3600$.
\end{exemple}

\section{8. Problèmes à étapes}
Une situation peut demander d’abord un volume, puis une masse à partir d’une masse volumique, puis une durée avec un débit.

\[
m=\rho V$ puis $t=V/q
\]

\begin{exemple}
Le raisonnement doit conserver les unités à chaque étape pour éviter une réponse numériquement correcte mais physiquement incohérente.
\end{exemple}

\section{Méthode générale de résolution}
\begin{methode}
\begin{enumerate}
\item Identifier chaque grandeur et son unité.
\item Écrire la relation avant de remplacer par les valeurs.
\item Convertir les unités avant l’opération si elles ne sont pas compatibles.
\item Utiliser les puissances de dix pour justifier les conversions d’aires ou volumes.
\item Conserver les unités dans les calculs intermédiaires.
\item Contrôler la plausibilité du résultat et préciser l’unité finale.
\end{enumerate}
\end{methode}

\section{Problème guidé et stratégie}
Dans un problème de troisième, la difficulté ne vient pas seulement du calcul. Il faut reconnaître la notion utile, choisir une représentation, enchaîner plusieurs étapes et expliquer pourquoi la méthode est valide. On commence donc par isoler les données, puis on écrit la relation mathématique avant de remplacer par des nombres. Une réponse est considérée comme complète lorsqu’elle comporte une justification, un calcul lisible, un contrôle et une interprétation.

\begin{attention}
Une figure, un graphique, un tableau ou une calculatrice fournit des informations, mais ne remplace pas une justification lorsque le problème demande une preuve. Inversement, un calcul exact sans interprétation peut rester incomplet dans un contexte concret.
\end{attention}

\section{Contrôler un résultat}
Avant de valider une réponse, on vérifie le signe, l’ordre de grandeur, les unités et la compatibilité avec les hypothèses. On peut aussi refaire le calcul par une autre représentation : formule et graphique, valeur exacte et approximation, calcul direct et vérification dans l’énoncé. Cette habitude permet de repérer une erreur de saisie ou une mauvaise modélisation avant la conclusion.

\section{Erreurs fréquentes}
\begin{itemize}
\item Mélanger kilomètres et mètres ou heures et secondes.
\item Convertir une aire avec un facteur de longueur au lieu de son carré.
\item Convertir un volume sans cuber le facteur.
\item Confondre masse volumique et concentration.
\item Oublier l’unité d’une grandeur quotient.
\item Utiliser une formule sans vérifier la nature des grandeurs.
\end{itemize}

\section{À retenir}
\begin{aretenir}
\begin{itemize}
\item Une grandeur quotient compare une quantité à une autre.
\item Les unités font partie du raisonnement.
\item Les conversions de surfaces et volumes utilisent des facteurs au carré ou au cube.
\item Les puissances de dix rendent les conversions transparentes et contrôlables.
\end{itemize}
\end{aretenir}

\section{Réinvestissement : Masse volumique}
Pour réinvestir cette notion, on part d’une situation où plusieurs représentations sont possibles. La masse volumique relie masse et volume. L’élève doit expliciter son choix, effectuer le calcul et revenir au contexte. La valeur numérique dépend des unités choisies : $1\,\text{g/cm}^3=1000\,\text{kg/m}^3$. Le contrôle final consiste à vérifier que la réponse respecte les hypothèses et que son ordre de grandeur est compatible avec les données.

\section{Réinvestissement : Vitesse}
Pour réinvestir cette notion, on part d’une situation où plusieurs représentations sont possibles. La vitesse moyenne est une distance divisée par une durée. L’élève doit expliciter son choix, effectuer le calcul et revenir au contexte. Les unités de $d$ et $t$ déterminent l’unité de $v$ ; on ne mélange pas kilomètres et secondes sans conversion. Le contrôle final consiste à vérifier que la réponse respecte les hypothèses et que son ordre de grandeur est compatible avec les données.

\section{Réinvestissement : Vitesse en unités différentes}
Pour réinvestir cette notion, on part d’une situation où plusieurs représentations sont possibles. Pour convertir $90\,\text{km/h}$ en mètres par seconde, on convertit simultanément la distance et la durée. L’élève doit expliciter son choix, effectuer le calcul et revenir au contexte. Le facteur usuel $3{,}6$ peut être retrouvé par $1000/3600$. Le contrôle final consiste à vérifier que la réponse respecte les hypothèses et que son ordre de grandeur est compatible avec les données.

\section{Réinvestissement : Vitesse}
Pour réinvestir cette notion, on part d’une situation où plusieurs représentations sont possibles. La vitesse moyenne est une distance divisée par une durée. L’élève doit expliciter son choix, effectuer le calcul et revenir au contexte. Les unités de $d$ et $t$ déterminent l’unité de $v$ ; on ne mélange pas kilomètres et secondes sans conversion. Le contrôle final consiste à vérifier que la réponse respecte les hypothèses et que son ordre de grandeur est compatible avec les données.
